\section{The Symptomatology of Trust}
\label{sec:symptomatology}

\noindent
Before a condition is diagnosed it is presented, and a diagnosis that proceeds without first attending to the symptoms is a diagnosis of nothing. This section assembles the symptoms. It does not yet explain them; explanation is the work of \xref{sec:diagnosis}, and to anticipate it here would be to read the symptoms in the light of a thesis they are meant to motivate rather than to confirm. The discipline of the section is therefore descriptive: it sets out a range of observable, recognisable phenomena, drawn from the most intimate scale of human dealing to the most public, and asks only that they be looked at together, in the expectation that what governs them will appear, when it appears, as the explanation of a pattern already seen rather than as a hypothesis imposed upon scattered instances.

The phenomena, looked at together, exhibit a single and initially paradoxical shape. \emc{The apparatus by which trust is secured proliferates, and trust itself grows scarcer.} Wherever one looks, the instruments designed to produce or guarantee or certify reliability multiply\,---\,longer contracts, more elaborate vows, denser systems of credentialing, ever finer mechanisms of rating and verification\,---\,and the multiplication does not arrest the erosion of the thing the instruments exist to secure but accompanies it, and, in a way the section will only describe and not yet account for, seems bound up with it. The relation between the proliferation and the erosion is left, here, deliberately unstated. That there is such a relation, recurring across scales that have otherwise little in common, is the observation the section exists to establish.

The phenomena are presented across three scales, from the intimate outward, not because the intimate is a miniature of the public\,---\,whether it is or is not is a question for later\,---\,but because the recurrence of one shape across scales that differ in every other respect is itself the first thing to be seen.

\subsection{The intimate and familial scale}
\label{sub:sympt-intimate}

At the most intimate scale, the symptom appears first in the very document from which this paper set out. A betrothal is undertaken, and around it gathers an apparatus of solemnisation\,---\,the auspicious day, the silk, the gravity of the language, in other settings the ring, the announcement, the assembled witnesses\,---\,whose elaboration is in no evident proportion to the security it confers, since it confers, in the strict sense, none. People who undertake such a thing know, when they reflect, that nothing in the ceremony binds; and they elaborate the ceremony anyway, and feel the elaboration as weighty rather than as idle. This is already a symptom: a solemnising apparatus is invoked with full seriousness in the acknowledged absence of anything it could enforce.

The symptom recurs, in a sharper and more troubling form, in the marriage that may follow. The vows exchanged at weddings have grown, across recent generations and in many places, more articulate, more personal, more lavish in their profession of permanence, at the same time as the dissolution of the unions they open has grown more common rather than less. The observation is not offered here as a causal claim in either direction; it is offered as a juxtaposition that the symptomatology must record, that the form of the promise has become more elaborate as the persistence of what it promises has become less assured. The two curves, of the eloquence of the vow and of the fragility of the bond, have moved, over the relevant period, in the same direction, and a symptomatology notes this without yet asking what to make of it.

The symptom appears again, and most pressingly for the concerns of this paper, in the predicament of the one who must win the trust of a family not yet his own. Here the apparatus that elsewhere proliferates is precisely what is unavailable. The suitor has, by hypothesis, no track record with the family, no accumulated history by which he might be known, no third party they already trust who can vouch for him; he has only a sincerity that is, to him, entirely real, and that has, to them, no form they can read as evidence. The symptom at this scale is thus not the proliferation of the apparatus but the felt insufficiency of everything that is not it: the discovery that sincerity, however genuine, does not present itself as ground, and that the very people whose trust matters most are the people to whom one has least that they will count.

\subsection{The legal and judicial scale}
\label{sub:sympt-legal}

At the scale of law and commerce the same shape recurs with a documentary precision that makes it easy to measure. Contracts have grown longer. The instruments by which commercial parties bind one another have expanded, across the last century and with sharp acceleration in recent decades, into documents of a length and intricacy no party reads in full, drafted in anticipation of every contingency and every adversarial reading, and lengthening still; and the lengthening has not been accompanied by a growing ease of dealing or a growing confidence between contracting parties, but, by the testimony of those who deal under such instruments, by the opposite. The thicker the contract, the less, not the more, the parties feel secure; the elaboration of the terms is experienced less as the building of confidence than as its pre-emptive substitute, an arming against a breach already half expected.

Beside the contract stands the memorandum of understanding, and it is, for this paper, the most telling instrument at this scale precisely because it is the contract's opposite. The memorandum of understanding is, by design, non-binding: it records an intention, sets out the shape of a contemplated dealing, and explicitly withholds the force of an enforceable obligation. Such instruments proliferate. They are drafted with care, signed with ceremony, and treated by their signatories as significant, in the full and stated knowledge that they bind nothing. That parties expend such effort on the solemn execution of documents they know to be unenforceable is a symptom of exactly the kind this section collects: a grave investment in a form whose want of force is not concealed but declared.

And behind both contract and memorandum stands the court, in which the symptom takes its most consequential form. The machinery of enforcement, the sword to which all these instruments are nominally referred, is under a strain that has begun to show. The volume of breach exceeds the capacity of the courts to adjudicate it; the cost of enforcement, in time and money and attention, rises until for many wrongs it exceeds the value of the remedy; the period over which a judgment is obtained lengthens, and the prospect of actually recovering on a judgment, once obtained, narrows. The sword, in a word, hangs but increasingly does not fall, and is known not to fall; and an instrument of enforcement that is known not to fall has ceased, whatever its formal availability, to function as the thing it was. The symptom here is not the proliferation of an apparatus but the hollowing of one: the persistence of the full formal machinery of compulsion alongside its growing practical inability to compel.

\subsection{The universal scale}
\label{sub:sympt-universal}

At the widest scale the shape appears in its purest and most general form, in the instruments by which strangers are made legible to strangers. The letter of recommendation is the clearest case. It is an instrument designed to transmit a knowledge that the recipient cannot acquire directly: the writer, who knows the candidate, vouches for the candidate to a committee, which trusts the writer. And it has, by common acknowledgment within the institutions that rely on it, very largely ceased to function. The letters have inflated until nearly every candidate is described as among the best the writer has encountered; the superlative, having become universal, has become uninformative; and committees, no longer able to extract a signal from a channel everyone has learned to saturate, turn away from the letters toward whatever can be counted\,---\,scores, citations, rankings\,---\,not because the counted things are richer but because they are, at least, harder to inflate. The apparatus of vouching has proliferated to the point of self-cancellation, and trust has retreated from it.

The same shape governs the systems of rating and certification by which the digital economy attempts to make its strangers legible. Trust badges, verification marks, star ratings, certification seals multiply across every surface on which strangers must deal with strangers; and their multiplication is met, by those who encounter them, not with growing confidence but with a growing and well-schooled wariness, since the very ubiquity of the mark of trustworthiness teaches that the mark is cheap, gameable, and as available to the predator as to the honest. The more the surface is covered with the signs of trustworthiness, the less any one of them reassures. Here the proliferation and the erosion are not merely concurrent but legibly bound: each new layer of certification is a response to the failure of the last, and itself elicits the next.

And beneath all these the most general symptom of the age, which the section records last because it underlies the rest: the condition of being surrounded by symbols whose provenance and whose intention cannot be established. The cost of producing a convincing token\,---\,a document, an image, a voice, a record of conduct\,---\,has fallen, with the technologies of the present, toward zero, and the consequence is a saturation of the human environment with signs that are no longer reliably tied to anything they purport to indicate. A person now moves through a world of symbols any of which may be confected, and the rational response to such a world\,---\,a generalised lowering of the weight given to all symbolic testimony\,---\,is itself a symptom of the first order: a withdrawal of trust not from this or that untrustworthy source but from the symbolic as such.

\subsection{A taxonomy of the symptoms}
\label{sub:sympt-taxonomy}

The phenomena gathered above are not a miscellany. Looked at together, they are so many appearances of a single fissure, opening at different scales and in different materials, between an apparatus that is meant to secure trust and the trust it fails to secure. To make the unity visible, and to give the diagnosis of \xref{sec:diagnosis} a structured object rather than a list of anecdotes to work upon, the symptoms may be set out along a small number of shared dimensions: the scale at which the symptom appears; the apparatus of assurance proper to that scale; the mode in which that apparatus fails; and the specific form the governing paradox\,---\,proliferation alongside erosion\,---\,takes in each case. The table that follows arranges them so. It is offered as a description and not as an explanation: each row records what is observed, and the column of failure-modes names the manner of a breakdown without yet assigning its cause. That every row instances one fissure, and what that fissure is, is the matter the table hands forward to the diagnosis.

\begin{center}
\small
\renewcommand{\arraystretch}{1.35}
\begin{longtable}{>{\RaggedRight\arraybackslash}p{2.7cm} >{\RaggedRight\arraybackslash}p{2.2cm} >{\RaggedRight\arraybackslash}p{3.2cm} >{\RaggedRight\arraybackslash}p{3.4cm}}
\toprule
\textcolor{rosedeep}{\textbf{\small Phenomenon}} & \textcolor{rosedeep}{\textbf{\small Scale}} & \textcolor{rosedeep}{\textbf{\small Apparatus of assurance}} & \textcolor{rosedeep}{\textbf{\small Mode of failure}} \\
\midrule
\endfirsthead
\toprule
\textcolor{rosedeep}{\textbf{\small Phenomenon}} & \textcolor{rosedeep}{\textbf{\small Scale}} & \textcolor{rosedeep}{\textbf{\small Apparatus of assurance}} & \textcolor{rosedeep}{\textbf{\small Mode of failure}} \\
\midrule
\endhead
\bottomrule
\endfoot
Betrothal, wedding vow & intimate & ceremony, solemnity, profession of permanence & elaboration grows as the bond's persistence weakens; gravity without enforceability \\
Winning a family's trust & familial & sincerity, profession, proposed conduct & sincerity does not present as evidence to those who ask for grounds \\
Commercial contract & legal & enumerated terms, penalties, remedies & thickening of terms experienced as anxiety, not confidence; arming against expected breach \\
Memorandum of understanding & legal & solemn execution of a declared non-binding form & grave investment in an instrument known to bind nothing \\
The court, enforcement & judicial & compulsion, the threat of the sanction & the sword hangs but does not fall and is known not to; formal machinery hollowed of practical force \\
Letter of recommendation & institutional & a trusted writer's vouching along a chain & inflation to the superlative; the channel saturated to uninformativeness \\
Rating, certification, trust badge & platform & marks, seals, verification, star scores & ubiquity teaches cheapness; the mark reassures less the more it covers \\
The confected token & universal & symbolic testimony as such & zero-cost production severs sign from referent; rational withdrawal from the symbolic \\
\end{longtable}
\end{center}

\noindent
Read down the columns, the table says more than any row. The scale ranges from two persons to a whole society, and the apparatus changes utterly across that range, from red silk to a star rating; yet the column of failure-modes rhymes. In every case a form that is meant to carry assurance is found, in the end, not to carry it: the form persists, elaborates, even proliferates, while the assurance drains out of it. Whether at the scale of a vow or of a verification mark, the gap that opens is the same gap\,---\,between the apparatus and the trust\,---\,and it is the recurrence of this one gap across a range that shares nothing else that constitutes the central symptom of the age. The proliferation of the apparatus and the erosion of the trust are, the table suggests without yet showing, two faces of one process. Naming that process, and showing why the proliferation should accompany rather than arrest the erosion, is the work to which the symptomatology now hands over.
